\ifdefined\gkver\else\def\gkver{student}\fi
\documentclass[\gkver]{gaokaozhenti}
\begin{document}

\chapter{2026年江苏}

\begin{enumerate}
\item
%% number: 1
%% paperName: 2026年江苏高考第 1 题 4 分
%% typeId: 1
%% type: 单选题
%% chapter: 磁场
%% point: 带电粒子在磁场中的运动
%% method: 无
%% score: 4
%% degree: 850
%% duplicateId: 0
%% body:
平面内存在垂直平面向外的匀强磁场，一电子垂直磁场射入，不计空气阻力。如图所示，则电子偏转的轨迹可能是\xzanswer{A}
\onepicture[w=4.8cm]{figs/01.png}
\fourchoices[answer=A]
{a}
{b}
{c}
{d}
%% answer: A
%% memo:
\memoanswer{
由于不计空气阻力，则电子做匀速圆周运动，且根据左手定则可知电子向左偏转，可知电子偏转的轨迹为圆弧的一部分。

故选 A。
}

\item
%% number: 2
%% paperName: 2026年江苏高考第 2 题 4 分
%% typeId: 1
%% type: 单选题
%% chapter: 电磁感应
%% point: 远距离输电
%% method: 无
%% score: 4
%% degree: 850
%% duplicateId: 0
%% body:
使用高压输电，若发电站的输出功率不变，输电线的总电阻不变，将输电线中的电流降为原来的一半，则输电线上损耗的功率变为原来的\xzanswer{B}
\fourchoices[answer=B]
{2 倍}
{$\frac{1}{4}$}
{不变}
{4 倍}
%% answer: B
%% memo:
\memoanswer{
由题知输电电流降为原来的 $\frac{1}{2}$，即 $I' = \frac{1}{2} I$

设输电线总电阻为 $R$，则有 $P_{\text{损}} = I^{2} R$，$P_{\text{损}}' = (I')^{2} R = \frac{1}{4} P_{\text{损}}$

即损耗功率变为原来的 $\frac{1}{4}$。

故选 B。
}

\item
%% number: 3
%% paperName: 2026年江苏高考第 3 题 4 分
%% typeId: 1
%% type: 单选题
%% chapter: 原子和原子核
%% point: 玻尔模型
%% method: 无
%% score: 4
%% degree: 850
%% duplicateId: 0
%% body:
氢原子基态轨道的能量为 $E_{1}$，激发态轨道的能量为 $E_{2}$，已知普朗克常量为 $h$。则电子从该激发态跃迁至基态轨道释放的光子频率为\xzanswer{D}
\fourchoices[answer=D]
{$E_{1}$}
{$E_{2}$}
{$\frac{E_{1} - E_{2}}{h}$}
{$\frac{E_{2} - E_{1}}{h}$}
%% answer: D
%% memo:
\memoanswer{
根据玻尔跃迁理论，氢原子从高能级（激发态）跃迁至低能级（基态）时，释放光子的能量等于两能级的能量差，即 $\Delta E = E_{2} - E_{1}$

再结合光子能量与频率的关系 $\Delta E = h\nu$

联立可得光子频率 $\nu = \frac{E_{2} - E_{1}}{h}$

故选 D。
}

\item
%% number: 4
%% paperName: 2026年江苏高考第 4 题 4 分
%% typeId: 1
%% type: 单选题
%% chapter: 曲线运动
%% point: 平抛运动
%% method: 无
%% score: 4
%% degree: 750
%% duplicateId: 0
%% body:
如图所示，小明跳起从同一高度将排球水平击出三次，排球分别落于 a、b、c 三点，不计空气阻力，关于该过程，下列说法正确的是\xzanswer{C}
\onepicture[w=6.5cm]{\includesvg{figs/04}}
\fourchoices[answer=C]
{排球的飞行时间满足 $t_{a} > t_{b} > t_{c}$}
{排球的飞行时间满足 $t_{a} < t_{b} < t_{c}$}
{排球的初速度满足 $v_{a} < v_{b} < v_{c}$}
{排球的初速度满足 $v_{a} = v_{b} = v_{c}$}
%% answer: C
%% memo:
\memoanswer{
AB．竖直方向上，根据自由落体运动公式 $h=\frac{1}{2}gt^{2}$，可得飞行时间 $t=\sqrt{\frac{2h}{g}}$

因为排球从同一高度抛出，且落在同一水平面上，所以下落高度 $h$ 相同，故三次飞行的时间相等，即 $t_{a}=t_{b}=t_{c}$，故 A 错误，B 错误；

CD．水平方向上，根据匀速直线运动公式 $x = v_{0}t$，可得初速度 $v_{0} = \frac{x}{t}$

由图可知，排球的水平位移关系为 $x_{a} < x_{b} < x_{c}$，且已知飞行时间 $t$ 相等，所以初速度满足 $v_{a} < v_{b} < v_{c}$，故 C 正确，D 错误。

故选 C。
}

\item
%% number: 5
%% paperName: 2026年江苏高考第 5 题 4 分
%% typeId: 1
%% type: 单选题
%% chapter: 原子和原子核
%% point: 半衰期
%% method: 无
%% score: 4
%% degree: 850
%% duplicateId: 0
%% body:
已知某种元素的半衰期为 39 万年，一群该元素对应的粒子的原始质量为 $m_{0}$，经历 12 万年后该粒子的质量剩余 $m_{1}$，则下列关于剩余的该粒子质量随时间变化的图像正确的是\xzanswer{A}
\fourchoices[answer=A, ispicture=true, h=2.7cm]
{figs/05a.png}
{figs/05b.png}
{figs/05c.png}
{figs/05d.png}
%% answer: A
%% memo:
\memoanswer{
根据半衰期公式 $m_{\text{余}}=m_{0}\left(\frac{1}{2}\right)^{\frac{t}{T}}$

这是一个以 $\frac{1}{2}$ 为底的指数函数，图像是一条单调递减、下凸的曲线。且由题知元素的半衰期 $T = 39$ 万年，则 39 万年时 $m_{\text{余}} = \frac{1}{2} m_{0}$

故选 A。
}

\item
%% number: 6
%% paperName: 2026年江苏高考第 6 题 4 分
%% typeId: 1
%% type: 单选题
%% chapter: 电场
%% point: 电场线
%% method: 无
%% score: 4
%% degree: 650
%% duplicateId: 0
%% body:
一带正电粒子沿直线穿越如图等势面，则下列说法正确的是\xzanswer{B}
\onepicture[w=5.0cm]{figs/06.png}
\fourchoices[answer=B]
{粒子电势能先增大后减小}
{粒子电势能先减小后增大}
{粒子所受电场力先做负功后做正功}
{粒子所受电场力不变}
%% answer: B
%% memo:
\memoanswer{
AB．由图可知，等势面越靠近中心电势越低，粒子沿直线从右侧外部向左侧运动，过程中电势先减小，后增大，粒子带正电，因此电势能先减小后增大，故 A 错误，B 正确；

C．电势能先减小后增大，说明电场力先做正功后做负功，故 C 错误；

D．等差等势面的疏密反映电场强度大小，越靠近中心等势面越密，电场强度越大，可知粒子所受电场力是变化的，D 错误。

故选 B。
}

\item
%% number: 7
%% paperName: 2026年江苏高考第 7 题 4 分
%% typeId: 1
%% type: 单选题
%% chapter: 曲线运动
%% point: 万有引力定律
%% method: 无
%% score: 4
%% degree: 750
%% duplicateId: 0
%% body:
发射一颗探月卫星，其先后在两个轨道上分别经过 M、N 两点，如图所示，关于卫星在 M、N 两点的情况，下列说法正确的是\xzanswer{D}
\onepicture[w=5.0cm]{figs/07.png}
\fourchoices[answer=D]
{在 N 点所受地球引力比月球引力大}
{在 N 点所受地球引力和月球引力相等}
{在 M 点所受月球引力比地球引力大}
{在 M 点所受地球引力比月球引力大}
%% answer: D
%% memo:
\memoanswer{
根据万有引力公式 $F = G \frac{Mm}{r^{2}}$，结合位置和轨道性质分析：

M 点分析：M 点距离地球更近、距离月球更远，且地球质量远大于月球质量，因此卫星在 M 点受到的地球引力大于月球引力；

N 点分析：N 点距离月球更近、距离地球更远，轨道向月球一侧弯曲，说明合力指向月球，即卫星在 N 点受到的月球引力大于地球引力。

故选 D。
}

\item
%% number: 8
%% paperName: 2026年江苏高考第 8 题 4 分
%% typeId: 1
%% type: 单选题
%% chapter: 光学
%% point: 光的干涉
%% method: 无
%% score: 4
%% degree: 750
%% duplicateId: 0
%% body:
如图所示为在重力作用下肥皂膜上出现的条纹。关于该条纹，下列说法正确的是\xzanswer{A}
\onepicture[w=4.2cm]{figs/08.jpg}
\fourchoices[answer=A]
{是干涉现象，自上而下相邻条纹间距越来越小}
{是干涉现象，自上而下相邻条纹间距相同}
{是衍射现象，自上而下相邻条纹间距越来越小}
{是衍射现象，自上而下相邻条纹间距相同}
%% answer: A
%% memo:
\memoanswer{
肥皂膜上的条纹是薄膜干涉现象，由肥皂膜前后两个表面的反射光叠加干涉形成。

竖直放置的肥皂膜受重力作用，液体向下流动，形成上薄下厚的结构：越靠下，肥皂膜厚度随竖直高度的增加越快（单位竖直距离的厚度变化量越大）。

根据薄膜干涉规律，相邻亮纹（或暗纹）对应的厚度差固定为 $\frac{\lambda}{2}$（$\lambda$ 为入射光波长），厚度变化率越大，相邻条纹的竖直间距越小，因此自上而下（从上到下）相邻条纹间距越来越小。

故选 A。
}

\item
%% number: 9
%% paperName: 2026年江苏高考第 9 题 4 分
%% typeId: 1
%% type: 单选题
%% chapter: 光学
%% point: 光的偏振
%% method: 无
%% score: 4
%% degree: 850
%% duplicateId: 0
%% body:
如图所示，由照相机拍摄的两组照片，甲图无法看清车内景物，乙图调整镜头后拍摄可以看清车内景物。关于该现象，下列说法正确的是\xzanswer{C}
\twopicture[numA=甲, numB=乙, labelprefix={}, h=3.0cm]
{figs/09a.png}
{figs/09b.png}
\fourchoices[answer=C]
{照相机镜头前加的是滤光片}
{照相机镜头前加的是增透膜}
{照相机镜头前加的是偏振片}
{照相机镜头前加的是增反膜}
%% answer: C
%% memo:
\memoanswer{
该现象的原理是：汽车玻璃表面的反射光是偏振光，甲图中玻璃反射的强光（眩光）盖过了车内景物的光线，因此无法看清车内景物。

A．滤光片仅用于过滤特定波长的光，无法消除玻璃的反射眩光，A 错误。

B．增透膜是镜头本身用来增加整体透光率的薄膜，不能消除玻璃表面的反射眩光，B 错误。

C．偏振片可以阻挡玻璃反射的偏振光，滤除反射眩光后，就能看清车内景物，符合题意，C 正确。

D．增反膜的作用是增强反射，更无法看清车内，D 错误。

故选 C。
}

\item
%% number: 10
%% paperName: 2026年江苏高考第 10 题 4 分
%% typeId: 1
%% type: 单选题
%% chapter: 电磁感应
%% point: 法拉第电磁感应定律
%% method: 无
%% score: 4
%% degree: 750
%% duplicateId: 0
%% body:
旅客在网上购买高铁票后，可使用二代身份证在自动检票闸机上刷身份证进出站。当身份证进入刷卡机感应范围后，刷卡机会向身份证发射高频电磁波，下列说法正确的是\xzanswer{B}
\fourchoices[answer=B]
{身份证放在刷卡机上不动不会产生感应电流}
{身份证放在刷卡机上不动就会产生感应电流}
{身份证快速靠近刷卡机感应电流会变大}
{身份证快速远离刷卡机感应电流会变大}
%% answer: B
%% memo:
\memoanswer{
感应电流产生的条件是闭合回路中磁通量发生变化，刷卡机发射的高频电磁波本质是周期性变化的磁场，结合法拉第电磁感应定律分析各选项。

A．身份证放在刷卡机上不动时，刷卡机产生的交变磁场会使穿过身份证内置线圈的磁通量随时间周期性变化，因此会产生感应电流，故 A 错误；

B．由 A 的分析可知，身份证不动时磁通量持续变化，会产生感应电流，故 B 正确；

C．根据法拉第电磁感应定律可知感应电流大小与磁通量变化率成正比，身份证快速靠近或者快速远离刷卡机过程，由于刷卡机产生的是交变磁场，这两个过程中的磁通量变化率都有可能是减小的，感应电流也可能是减小的，故 CD 错误。

故选 B。
}

\item
%% number: 11
%% paperName: 2026年江苏高考第 11 题 4 分
%% typeId: 1
%% type: 单选题
%% chapter: 运动和力
%% point: 动量定理
%% method: 无
%% score: 4
%% degree: 450
%% duplicateId: 0
%% body:
如图所示，水平面内一光滑小球沿正方形线框 $fghi$ 保持速率不变运行，$a$、$b$、$c$、$d$ 分别是各边中点，在拐点处小球所受合力的冲量大小和其动量大小的比值是\xzanswer{A}
\onepicture[w=5.5cm]{\includesvg{figs/11}}
\fourchoices[answer=A]
{$\frac{ab}{gb}$}
{$\frac{ab}{bd}$}
{$\frac{ab}{fh}$}
{$\frac{ab}{hi}$}
%% answer: A
%% memo:
\memoanswer{
根据动量定理，小球在拐点处受到合力的冲量等于动量的变化量，即 $I = \Delta p$

设小球速率为 $v$，动量大小 $p = mv$

正方形相邻边方向垂直，小球拐弯前后速度大小均为 $v$，方向夹角为 $90^{\circ}$，因此动量变化量的大小为：$|\Delta p| = \sqrt{(mv)^{2} + (mv)^{2}} = \sqrt{2}mv = \sqrt{2}p$

可得 $\frac{I}{p} = \sqrt{2}$

结合几何关系，选项中只有 A 符合。

故选 A。
}

\item
%% number: 12
%% paperName: 2026年江苏高考第 12 题 15 分
%% typeId: 4
%% type: 实验题
%% chapter: 电场
%% point: 电容器充放电实验
%% method: 无
%% score: 15
%% degree: 550
%% duplicateId: 0
%% body:
某实验小组用手机记录电容器的充放电过程。
\begin{enumerate}
	\item
	如图~\ref{2026江苏12a} 所示，电容器的电容为 \tkanswer[2.4cm]{$1\times10^{-3}\UF$}。
	\item
	关于电容器的电容，下列说法正确的是 \tkanswer[1cm]{C}。

	\threechoices[answer=C]
	{随着放电的进行，电容器的电容会逐渐变小}
	{随着充电的进行，电容器的电容会逐渐变小}
	{无论是否充放电，电容器的电容均不变}
	\item
	用手机记录电容器充满电后的放电过程，电路图如图~\ref{2026江苏12b} 所示，则下列步骤排序应为 \tkanswer[1.4cm]{②①③}。

	\begin{steps}
		\item
		S 接 b 至 $\mathrm{A}_{2}$ 为 0
		\item
		打开手机开始拍摄，S 接 a 至 $\mathrm{A}_{1}$ 为 0
		\item
		关闭手机
	\end{steps}
	\item
	在图~\ref{2026江苏12c} 中定性描绘出电容器放电过程电路中的电流随时间变化的规律。
	\item
	已知电磁式电流表中有铝框，上面绕有线圈，请简述这一结构的作用：\tkanswer[10cm]{铝框转动时会产生感应电流，感应电流的安培力产生电磁阻尼，阻碍指针摆动，使指针快速稳定，便于读数}。
\end{enumerate}
\threepicture[labelA=2026江苏12a, labelB=2026江苏12b, labelC=2026江苏12c, h=3.2cm, align=right]
{figs/12a.png}
{figs/12b.png}
{\includesvg{figs/12c}}
%% answer:
\jdanswer{
\begin{enumerate}
	\item
	$1\times10^{-3}$

	\item
	C

	\item
	②①③

	\item

	\item
	铝框转动时会产生感应电流，感应电流的安培力产生电磁阻尼，阻碍指针摆动，使指针快速稳定，便于读数

\end{enumerate}
}
%% memo:
\memoanswer{
（1）根据题图可知电容器的电容为 $1000\UmF=1\times10^{-3}\UF$。

（2）电容是电容器本身的固有属性，由电容器自身结构（极板正对面积、极板间距、电介质）决定，与是否充放电、带电量、极板间电压无关，因此无论是否充放电，电容均不变。故选 C。

（3）实验目的是记录充满电后的放电过程，步骤为：先对电容器充电，开关接 a 到充电回路，充电至电流为 0（电容充满电），打开手机准备记录，再将开关切换到 b 接入放电回路开始放电，放电至电流为 0（电容放完电），放电完成后关闭手机，对应步骤排序为②①③（先充电、再放电记录、最后关机）。

（4）电容器放电时，初始电流最大，随时间推移，极板电荷量减少，电压降低，电流逐渐减小，最终电流为 0，且电流减小速率越来越慢，如图所示。
\onepicture[w=4.5cm]{\includesvg{figs/12_an}}

（5）铝框是导体，指针摆动带动铝框转动，切割磁感线产生感应电流，根据楞次定律，安培力会阻碍铝框的相对运动，即电磁阻尼，能让指针快速停止摆动，方便读数。
}

\item
%% number: 13
%% paperName: 2026年江苏高考第 13 题 6 分
%% typeId: 6
%% type: 计算题
%% chapter: 气体性质
%% point: 热力学第一定律
%% method: 无
%% score: 6
%% degree: 750
%% duplicateId: 0
%% body:
一定质量的理想气体从 a 状态到 b 状态经历等温变化，此过程外界压缩气体对其做的功为 $W$。已知 a 状态气体体积为 $V_{0}$，压强为 $p_{0}$，b 状态气体体积为 $V$。求：
\begin{enumerate}
	\item
	该过程气体放出的热量；
	\item
	b 状态时气体的压强。
\end{enumerate}
%% answer:
\jdanswer{
\begin{enumerate}
	\item
	$W$

	\item
	$p = \frac{V_{0}}{V}p_{0}$

\end{enumerate}
}
%% memo:
\memoanswer{
（1）理想气体经历从 a 状态到 b 状态等温变化，气体的内能不变，由热力学第一定律有 $\Delta U = Q + W = 0$

解得 $Q = -W$

负号表示过程中气体放出热量 $W$。

（2）理想气体经历从 a 状态到 b 状态等温变化，由玻意耳定律有 $p_{0}V_{0} = pV$

解得 $p = \frac{V_{0}}{V}p_{0}$
}

\item
%% number: 14
%% paperName: 2026年江苏高考第 14 题 8 分
%% typeId: 6
%% type: 计算题
%% chapter: 曲线运动
%% point: 竖直平面圆周运动
%% method: 无
%% score: 8
%% degree: 550
%% duplicateId: 0
%% body:
一质量 $m = 2\times10^{3}\Ukg$ 的汽车经过一半径 $R = 20\Um$ 的拱形桥顶部时，速度 $v = 10\Ums$。已知牵引力 $F = 2000\UN$，阻力 $f = 1000\UN$。求此时：
\begin{enumerate}
	\item
	汽车的向心加速度大小；
	\item
	汽车受到的合力大小。
\end{enumerate}
%% answer:
\jdanswer{
\begin{enumerate}
	\item
	$5\Umsq$

	\item
	$1000\sqrt{101}\UN$

\end{enumerate}
}
%% memo:
\memoanswer{
（1）根据题意，由向心加速度公式可得 $a_{n}=\frac{v^{2}}{R}=\frac{10^{2}}{20}\Umsq=5\Umsq$。

（2）汽车在最高点时，水平方向上有 $F_{x\text{合}} = F - f = 1 \times 10^{3}\UN$

竖直方向有 $F_{y\text{合}} = ma_{n} = 1 \times 10^{4}\UN$

则汽车在最高点受到的合力大小为 $F_{\text{合}} = \sqrt{F_{x\text{合}}^{2} + F_{y\text{合}}^{2}} = 1000 \sqrt{101}\UN$
}

\item
%% number: 15
%% paperName: 2026年江苏高考第 15 题 12 分
%% typeId: 6
%% type: 计算题
%% chapter: 电磁感应
%% point: 电磁感应受力分析
%% method: 无
%% score: 12
%% degree: 450
%% duplicateId: 0
%% body:
扫地机器人可利用磁场对通电线圈的安培力实现悬浮与驱动。如图，一扫地机器人内含 $abcd$、$a'b'c'd'$ 两矩形线框，其中 $a'b'c'd'$ 为平衡框，用于保持悬浮，$abcd$ 为驱动框，用于提供向前的动力。上下两导线中通电后，其间形成的磁场可视为匀强磁场，磁感应强度大小为 $B$，方向垂直纸面向里。每个框都有一条边长为 $L$ 并嵌入金属芯，金属芯使得该边处磁场的磁感应强度变为原来的 $k$（$k>1$）倍。两线框通过绝缘杆固定连接，总质量为 $m$，运动过程中线框整体受到的摩擦阻力大小恒定，导线和线框均始终在同一竖直平面内，且线框的上下边与导线始终水平。（重力加速度为 $g$）
\begin{enumerate}
	\item
	为了使线框整体悬浮，求 $a'b'c'd'$ 框中的电流 $I$ 的大小和方向；
	\item
	若线框向右运行速度大小为 $v_{1}$，求 $ab$ 边的感应电动势 $E_{1}$；
	\item
	若线框以速度大小 $v_{0}$ 在水平方向上匀速运动，某时刻保持电流大小不变，仅将动力线框 $abcd$ 中的电流反向，线框在安培力和阻力作用下做匀减速直线运动，经过位移 $s$ 后停止，求整个减速过程，动力线框 $cd$ 边所受安培力的平均功率 $P$。
\end{enumerate}
\onepicture[w=9.0cm, align=right]{\includesvg{figs/15}}
%% answer:
\jdanswer{
\begin{enumerate}
	\item
	$I = \frac{mg}{(k - 1)BL}$，电流方向为逆时针

	\item
	$E_{1} = BLv_{1}$

	\item
	$\frac{kmv_0^3}{8s(k - 1)}$

\end{enumerate}
}
%% memo:
\memoanswer{
（1）根据题意可知，为了使线框整体悬浮，竖直方向上有 $kBIL - BIL = mg$

解得 $I = \frac{mg}{(k - 1)BL}$

由左手定则可知，电流方向为逆时针。

（2）若机器人向右运行速度为 $v_{1}$，$ab$ 边的感应电动势 $E_{1}=BLv_{1}$

（3）根据题意，若某时刻线框的水平速度为 $v_{0}$，由于 $cd$ 边所受安培力一直为 $ab$ 边所受安培力的 $k$ 倍，设 $cd$ 边所受安培力大小为 $F$，则 $ab$ 边所受安培力大小为 $\frac{F}{k}$，则有 $F - \frac{F}{k} = f$

保持电流大小不变，方向反向，则安培力方向反向，由牛顿第二定律有 $f + F - \frac{F}{k} = ma$

又有 $v_{0}^{2}=2as$

解得线框 $cd$ 边所受安培力大小为 $F=\frac{kmv_{0}^{2}}{4s\left(k-1\right)}$

线框 $cd$ 边所受安培力的平均功率 $P = F \bar{v} = F \cdot \frac{v_{0}}{2} = \frac{km v_{0}^{3}}{8s(k-1)}$
}

\item
%% number: 16
%% paperName: 2026年江苏高考第 16 题 15 分
%% typeId: 6
%% type: 计算题
%% chapter: 运动和力
%% point: 动量和能量
%% method: 无
%% score: 15
%% degree: 350
%% duplicateId: 0
%% body:
如图所示，在光滑水平面上固定两个柱形光滑轨道，轨道上分别约束着只能沿轨道方向运动的两个小球 B 和 C，质量均为 $\frac{3}{5}m$。小球 B 和 C 通过弹性限度足够大的相同轻质弹簧与质量为 $m$ 的小球 A 相连。初始时，两弹簧均处于原长。现有一质量为 $M$ 的小球 D 以速度 $v_{0}$ 沿轨道方向与小球 A 发生对心弹性碰撞，碰撞时间极短，不计空气阻力。
\begin{enumerate}
	\item
	求小球 D 与 A 碰撞后瞬间小球 A 的速度大小；
	\item
	若发生碰撞后小球 D 不再与 A 碰撞，求每根弹簧所具有的最大弹性势能；
	\item
	要使弹簧第一次恢复原长时，小球 D 与 A 恰好再次发生碰撞，求 $\frac{m}{M}$ 的值。
\end{enumerate}
\onepicture[w=8.0cm, align=right]{\includesvg{figs/16}}
%% answer:
\jdanswer{
\begin{enumerate}
	\item
	$v_{A} = \frac{2M}{M + m}v_{0}$

	\item
	$\frac{6mM^{2}v_{0}^{2}}{11(M + m)}$

	\item
	$\frac{1}{11}$

\end{enumerate}
}
%% memo:
\memoanswer{
（1）根据题意可知，小球 D 与小球 A 发生弹性正碰，由动量守恒定律有 $Mv_{0}=Mv_{1}+mv_{A}$

由能量守恒定律有 $\frac{1}{2}Mv_{0}^{2}=\frac{1}{2}Mv_{1}^{2}+\frac{1}{2}mv_{A}^{2}$

联立解得 $v_{1}=\frac{M-m}{M+m}v_{0}$，$v_{A}=\frac{2M}{M+m}v_{0}$

（2）根据题意可知，小球 A、B、C 共速时，弹簧的形变量最大，弹簧弹性势能最大，由动量守恒定律有 $mv_{A} = \left(m + 0.6m + 0.6m\right)v_{\text{共}}$

由能量守恒定律有 $\frac{1}{2}mv_{A}^{2}=\frac{1}{2}\left(m+0.6m+0.6m\right)v_{\text{共}}^{2}+2E_{p}$

解得 $E_{p}=\frac{3}{22}mv_{A}^{2}=\frac{6mM^{2}v_{0}^{2}}{11(M+m)}$

（3）根据题意，设小球 D 与小球 A 发生弹性正碰后到弹簧再次恢复原长的时间为 $t$，由动量守恒定律有 $mv_{A} = mv_{A}' + \left(0.6m + 0.6m\right)v'$

两边同时乘以 $t$ 并求和可得 $mv_{A}t = mx + \left(0.6m + 0.6m\right)x$

要使弹簧恢复原长时，$m$ 与 $M$ 能再次发生碰撞，则有 $x = v_{1}t$

整理可得 $\frac{2M}{M+m}v_{0}t=\frac{11}{5}\cdot\frac{M-m}{M+m}v_{0}t$

解得 $M = 11m$

即球质量 $m$ 和大球质量 $M$ 的质量比 $\frac{m}{M} = \frac{1}{11}$。
}

\end{enumerate}

\end{document}
